\documentclass{article}
\usepackage{amsmath,amssymb,hyperref}
\title{When a failed read spends the time needed for fallback}
\begin{document}
\maketitle

In the frozen two-report corridor case, a useful channel can save a tick. Trying a channel that may be ambiguous can also cost a tick: after waiting to listen, the best worst-case completion is nine ticks, whereas skipping the read gives eight. A selector that must preserve the eight-tick guarantee therefore needs evidence about the read before it spends the tick. Detecting a failed read afterward does not give that tick back.

\paragraph{Task, evidence, and costs.}
Keep the physical moves, report handoff, egress, and first delivery deadline of \href{https://github.com/eternal-november/eos-research-programmes/pull/372}{the optical rendezvous study}. At tick 2, B is at the hub after delivering its first report. It can leave immediately or spend the tick listening. The source A is at the refuge at this cut, attempts its cue, and then moves to its viable corridor. A's fixed action costs the same in both policies; A need not know whether B listened. The channel is available at this single cut. A read-start time, any receiver-specific coding, and the time at which a usable value becomes available are all indexed. No passive read while moving, second channel, or test that restores the spent tick is included.

There are two possible blocked branches for every admitted public pre-read code \(c\). For each, let \(S_0(c)\) and \(S_1(c)\) give all the outputs available to B by tick 3 if B listens. These sets are declared lookahead possibilities, not the outcome of this attempt. They include a symbol for no usable packet. Correctly ascertained absence may name a branch if it can occur in only one of the two cases. The source claim must bind the code, source action, branch, and read window; undecoded payload, delay or silence alone is no such claim.

\paragraph{Exact result.}
If B leaves the hub at tick 3 for the correct corridor, it can rendezvous at tick 4, perform the physical handoff during the next tick, and return to base by 7. If it takes the wrong corridor, it must travel back through the hub to the other one: the two extra moves put completion at 9. A cannot travel to the hub and return within its two-move allowance, and waiting there or visiting base cannot reveal the relevant blocked edge. Any local policy choosing between the corridors from an output in both supports must be wrong in at least one compatible history. Conversely, disjoint supports identify the branch. Thus listening has optimal worst-case completion 7 for disjoint supports and 9 for overlapping supports. These are guarantees over admitted histories, not a mixture-weighted service measure. Randomization cannot improve the all-histories guarantee.

Without listening, choosing a corridor at tick 2 and revising at the first wrong one guarantees 8, and the adverse branch needs 8. A remains at the correct meeting point independently of B's channel choice; if B arrives on the correct branch on its first try it can wait one tick for A. The exact pre-read rule for an eight-tick service contract is therefore to listen only in codes with disjoint supports and otherwise depart immediately. A success-qualified read supports seven, with no gain when contents or timing are still ambiguous. If additional evidence before departure has its own material cost, this rule needs that cost analyzed separately. An optional read that is free in the current model would be a different case.

\paragraph{A measured modality and what its metric forgets.}
\href{https://arxiv.org/html/2602.02249v1}{Putz et al.} compare acoustic schemes on real devices. Their total error rate (TER) averages decoded bit error and codes a failed decode as total error. In an authored four-bit control, let the first bit identify the branch and the others be nuisance bits, with both branches equally likely. In the first law, branch zero erases half its packets and otherwise sends every bit correctly; branch one always delivers its branch bit but gets two nuisance bits wrong. In a second law, each branch has probability one quarter of erasure, one half of one nuisance error, and one quarter of two nuisance errors. Both laws have aggregate erasure one quarter and per-branch TER one half. Absence identifies a branch in just the first law. This distinguishes the channel interfaces even before taking time into account. It does not treat partial payloads as valid application packets or estimate either law from the paper.

The published testbed decodes recordings offline after segmentation. Neither BER/TER nor the fact that an offline decoder can recover a packet certifies availability by the robot's decision cut. A source adapter would still need an issue-time anchor, a causal read window and decoder cost, and the set of valid branch outputs by that cut, including closed/absent cases. The authors' independent acoustic study helps specify this missing validation. It was not replayed as a substitute for a deadline-aware trace. No new physical tests were performed.

\paragraph{Replay and boundary.}
A directly enumerated first-choice search agrees with the route proof on all 49 nonempty support pairs over three read outcomes, with and without required listening, and on two mixed pre-read regimes. The physical continuations execute the pinned corridor model. Whenever the support excludes one branch, the adapter maps the unique inferred branch to the frozen receiver's expected code; the unchanged 11-entry local policy completes at tick 7 in all 30 such finite translations. On ambiguous outputs the adapter cannot provide that code. The direct calculation ties. This is a fallback and measurement boundary for R1, not independent non-radio mission transfer or closure.
\end{document}
